\begin{helper}
Fix \(s>0\).  Suppose that for some threshold \(k_0\) and some constant
\(C>0\), every \(k\ge k_0\) with \(k\ge2\) satisfies
\[
  C\frac{k^{s-2}}{(\log k)^{2s-6}}\le r(s,k).
\]
Then there is a constant \(c>0\) such that every \(k\ge2\) satisfies
\[
  c\frac{k^{s-2}}{(\log k)^{2s-6}}\le r(s,k).
\]
\end{helper}
\begin{proof}
For \(k\ge2\), set
\[
  T(k)=\frac{k^{s-2}}{(\log k)^{2s-6}}.
\]
Since \(k>1\), we have \(\log k>0\), and hence \(T(k)>0\).

Let \(A=\{k\in\mathbb N: 2\le k<k_0\}\).  If \(A\) is empty, then every
\(k\ge2\) also satisfies \(k\ge k_0\), so the assumed large-\(k\) bound holds
with \(c=C\).

Now suppose \(A\) is nonempty.  For each \(k\in A\), \(\noderef{RamseyNumberPositive}\)
gives \(r(s,k)>0\), and \(T(k)>0\), so the ratio \(r(s,k)/T(k)\) is positive.
Because \(A\) is finite and nonempty, this ratio has a minimum over \(A\); call
one minimizing value \(R_{\min}>0\).  Put
\[
  c=\min\{C,R_{\min}\}.
\]
Then \(c>0\).  If \(k\ge k_0\), the inequality with \(c\) follows from
\(c\le C\), \(T(k)>0\), and the assumed large-\(k\) inequality.  If
\(2\le k<k_0\), then \(k\in A\), so \(R_{\min}\le r(s,k)/T(k)\).  Since
\(c\le R_{\min}\) and \(T(k)>0\), multiplying by \(T(k)\) gives
\[
  cT(k)\le r(s,k).
\]
These two cases cover all \(k\ge2\), proving the claim.
\end{proof}
